\documentclass[12pt,a4paper]{amsart}

\input{preamble}

\title[Exercises template example]{Exercises template example}
\author{Author}
\date{\today}

\begin{document}

\maketitle

This example demonstrates the current exercise-sheet template: exercises,
subparts, displayed mathematics, references between exercises, and visible
solutions.

\begin{xca}{1}\label{exercise:basic}
Let \(V\) be a finite-dimensional vector space.
\begin{enumerate}[label=(\alph*)]
  \item \label{item:dual-dual-map} Define the natural map
  \(V \to V^{\vee\vee}\).
  \item Prove that this map is an isomorphism.
\end{enumerate}
\end{xca}

\begin{solution}
The natural map sends \(v \in V\) to the functional
\[
  \lambda \longmapsto \lambda(v)
\]
on \(V^\vee\). After choosing a basis, this map is identified with the usual
identification of a finite-dimensional vector space with its double dual.
\end{solution}

\begin{xca}{2}\label{exercise:exact-sequence}
Let
\[
  0 \to A \to B \to C \to 0
\]
be a short exact sequence of vector spaces. Use \Cref{exercise:basic} to
describe the dual exact sequence.
\end{xca}

\begin{solution}
Dualizing gives an exact sequence
\[
  0 \to C^\vee \to B^\vee \to A^\vee \to 0.
\]
The reference to \Cref{exercise:basic} is meant to check cross-references in
the exercise template.
\end{solution}

\end{document}
